\section{ULP: наблюдение, пока основные ядра спят}\label{les:ulp}
\lessonintro
  {собрать проект с программой CPU и программой сопроцессора}
  {уроки 9, 27 и 35}
  {пробуждение после десяти последовательных нажатых отсчётов}
  {кнопка GPIO4--GND, резистор 10 кОм GPIO4--3V3, ESP-IDF 5.4.x}

\subsection{Два домена исполнения}
Deep-sleep из урока 9 будил CPU по таймеру или внешнему сигналу. ULP может опрашивать
поддерживаемые RTC-входы и будить основную программу при выполнении условия.
У ESP32-S3 есть ULP FSM и ULP RISC-V; здесь выбираем RISC-V и язык C.
Это отдельный маломощный сопроцессор, а не третье ядро LX7. Он не запускает Arduino,
Wi-Fi и обычные задачи FreeRTOS. Набор доступной периферии проверяют отдельно.

\begin{figure}[H]\centering
\begin{tikzpicture}[node distance=10mm]
\node[blk,fill=blkper] (a) {Кнопка\\GPIO4};
\node[blk,fill=blkrtc,right=of a] (b) {ULP RISC-V\\опрос 100 мс};
\node[blk,fill=blkcpu,right=22mm of b] (c) {Основные ядра\\сон / обработка};
\draw[arr] (a)--(b);\draw[arr] (b)--node[above,font=\scriptsize,align=center]{10 отсчётов\\ LOW}(c);
\end{tikzpicture}\caption{Простой частый опрос отделён от сложного редкого действия}\end{figure}

Кнопка подключена к RTC-входу \pin{4}: внешний резистор \SI{10}{k\ohm} подтягивает вход к 3V3, кнопка
замыкает его на GND. Встроенную подтяжку для этого входа в коде отключают, поэтому внешний резистор обязателен.

\begin{figure}[H]
\centering
\begin{circuitikz}[scale=0.95,transform shape]
  \node[blk,fill=blkcpu,minimum width=2.8cm,minimum height=3.8cm] at (0,0) {};
  \node[font=\sffamily\small\bfseries] at (0,2.3) {ESP32-S3};
  \foreach \y/\n in {1.2/3V3,0/{\pin{4}},-1.2/GND}
    \node[font=\sffamily\scriptsize,anchor=east] at (1.4,\y) {\n};
  \draw[thick,wred]    (1.4,1.2) -- (3.2,1.2);
  \draw[thick,wyellow] (1.4,0) -- (5.6,0);
  \draw[thick,wblack]  (1.4,-1.2) -- (5.6,-1.2);
  \draw[thick] (3.2,1.2) to[R,l_=\SI{10}{k\ohm}] (3.2,0);
  \draw[thick] (5.6,0) to[push button,l_=S1] (5.6,-1.2);
  \fill[wyellow] (3.2,0) circle (1.6pt);
  \fill[wyellow] (5.6,0) circle (1.6pt);
  \node[font=\sffamily\scriptsize,text=accent] at (4.4,0.35) {RTC GPIO};
\end{circuitikz}
\caption{Подключение кнопки к \pin{4}: внешняя подтяжка \SI{10}{k\ohm} к 3V3 и кнопка на GND.
Светодиод с этого вывода нужно снять}
\end{figure}

\subsection{Проект с двумя образами}
Основной код компилируется для Xtensa, ULP --- для RISC-V. CMake встраивает маленький образ
в основную прошивку и создаёт заголовок общих символов.
Полный проект \code{examples/ulp\_button} содержит корневой CMakeLists,
\code{main/CMakeLists.txt}, \code{main/main.c}, \code{main/ulp/main.c} и sdkconfig.defaults.
В defaults включены ULP RISC-V и резерв RTC-памяти. Открывать эти файлы как скетч .ino нельзя.
Сборка: \code{idf.py set-target esp32s3}, \code{idf.py build},
\code{idf.py -p COM5 flash monitor}; замените порт своим.

\subsection{Что делает основная программа}
GPIO4 настраиваем как RTC-вход, LED с него снимаем. Внешний резистор задаёт HIGH,
кнопка замыкает на GND. Программа ждёт секунду для наблюдения лога, загружает ULP,
задаёт период 100000 мкс, разрешает ULP-пробуждение и засыпает.

\codecap{Основные ядра: загрузить ULP и уснуть}
\codeinput{l36_main}

\subsection{Программа сопроцессора}
\codecap{ULP: десять последовательных отсчётов LOW}
\codeinput{l36_ulp}

Счётчик в RTC slow memory сохраняется между запусками ULP по таймеру. HIGH сбрасывает его,
LOW увеличивает; на десятом отсчёте запрашиваем пробуждение. Выход из main завершает
текущий опрос до следующего периода. Удержание длится примерно секунду, но результат зависит
от фазы опроса: это не точное измерение времени. После пробуждения отпустите кнопку
в течение секундной паузы, иначе следующий цикл снова обнаружит нажатие.

\subsection{Проверка экономии}
ULP полезен, если предотвращает достаточно дорогие пробуждения CPU.
Сравните активный опрос, пробуждение CPU каждые 100 мс и ULP с редким wakeup.
Измеряйте ток и длительность каждого режима. На DevKit стабилизатор, индикатор питания
и USB-мост могут перекрыть выгоду чипа. Для оценки автономности нужен ток всей целевой платы,
включая датчики, а не только значение deep-sleep из datasheet.

\begin{trouble}
Не создан заголовок ULP: проверьте CMake и конфигурацию. Нет wakeup: не тот RTC-вход
или не разрешён источник. Постоянное пробуждение: удерживается кнопка либо плавает вход.
Полные исходники приведены, но сборка ESP-IDF и измерения на плате в этой среде не выполнялись.
\end{trouble}
\begin{summary}
CPU и ULP имеют отдельные программы и ресурсы. Начинайте с условия пробуждения,
периода опроса и измеримого бюджета энергии.
\end{summary}

\begin{selfcheck}
\begin{enumerate}
\item Почему ULP не собирается обычным Arduino-скетчем?
\item Равны ли десять отсчётов точной секунде?
\item Почему ULP почти не меняет ток некоторых DevKit?
\end{enumerate}
\end{selfcheck}
\needspace{20\baselineskip}
\begin{tasks}
\begin{enumerate}
\item \lv{1} Выделите на схеме работающие во сне блоки.
\item \lv{1} Соберите проект и подтвердите источник wakeup в логе.
\item \lv{2} Сравните пороги 10 и 5 при одном периоде.
\item \lv{2} Добавьте общий счётчик опросов и прочитайте после wakeup.
\item \lv{3} Сравните заряд за час в трёх режимах при одинаковых событиях.
\item \lv{4} Замените кнопку поддерживаемым ULP-датчиком и добавьте порог с гистерезисом.
\end{enumerate}
\end{tasks}

\paragraph{Документация для продолжения}\href{https://docs.espressif.com/projects/esp-idf/en/v5.4/esp32s3/api-reference/system/ulp-risc-v.html}{Официальное руководство по теме}.
